\section{Conclusion: Design Justification and Security Guarantees}

This report presents a comprehensive secure data hosting server based on the Myers-Liskov Decentralized Model for Information Flow Control. The design enables hospitals and medical institutions to offer data sharing, multi-stakeholder collaboration, and untrusted program execution while maintaining formal confidentiality, integrity, and non-interference guarantees. This section justifies all design choices against project guidelines, summarizes security achievements, and compares the approach to alternative information flow frameworks.

\subsection{Design Choices Justification}

The project guidelines (Section 1) offer multiple design dimensions. We made explicit choices for each:

\subsubsection{Lattice Model: Fixed vs. Customizable}

\textbf{Our Choice}: Customizable per-principal lattices within decentralized model.

\textbf{Why}: Hospital has hundreds of patients and institutional roles (doctors, nurses, researchers, administrators). Fixed 3-level lattice (Low, Medium, High) cannot capture nuanced policies. Patient wants data readable by personal doctor only, not all hospital staff. Researcher wants data readable by study team only, not hospital administrators.

\textbf{Decentralized Solution}: Each principal defines own label independently ($\text{patient}: \{\text{readers}\}$). Patient specifies doctor; researcher specifies study team. Labels compose via join (intersection of readers). Scales to arbitrary principals with heterogeneous policies.

\textbf{Alternative Rejected}: Centralized global lattice (Volpano style). Simpler but inflexible. Forces all principals into pre-defined security levels. Doesn't support hospital scenario where each patient has independent policies.

\subsubsection{Declassification: Yes or No}

\textbf{Our Choice}: Yes, with owner-authorized declassification.

\textbf{Why}: Medical research requires data sharing workflows. Patient agrees to research use → de-identification → public aggregates. Three-step process requires controlled downgrading at each stage. Without declassification, once data labeled private, cannot ever be released without data destruction.

\textbf{Why Not Just Encryption}: Encryption without information flow analysis cannot distinguish between authorized and unauthorized decryption. Attacker with key can decrypt secret data; server cannot guarantee key holder is authorized. Information flow analysis (type-checking) adds layer: only authorized computation paths can see secrets.

\textbf{Formal Justification}: Owner-authorized declassification is legitimate (not violation) when justified by formal condition (k-anonymity, aggregation, etc.). All declassifications audited. Enables research while maintaining accountability.

\textbf{Alternative Rejected}: No declassification. Data labeled private stays private. Simpler security model but prevents medical research, data sharing, collaboration. Unrealistic for hospitals with research missions.

\subsubsection{Program Model: Fixed vs. User-Uploaded vs. Hybrid}

\textbf{Our Choice}: Hybrid white-box + black-box support.

\textbf{Why}: Hospital has two categories of programs:
\begin{itemize}
\item \textbf{Research programs}: Custom machine learning models, image analysis algorithms written by researchers. Source code available. Can type-check for information flow properties.
\item \textbf{Legacy programs}: Pre-existing image filters, data transformation tools. No source code available (vendor binaries). Cannot type-check; must use conservative labeling.
\end{itemize}

Hybrid approach enables both: researchers write white-box programs for maximum security assurance; hospital uses legacy black-box programs with conservative labeling. Both guaranteed secure by design.

\textbf{White-Box Design}: Type-checker verifies programs before execution. Three rules enforced:
\begin{enumerate}
\item Explicit flow: labels respect ordering ($L_1 \leq L_2$)
\item Implicit flow: program counter tracks control-flow labels
\item Declassification: only owner principals can declassify
\end{enumerate}

Type-check pass = program certified secure. No runtime enforcement needed; static verification sufficient.

\textbf{Black-Box Design}: Binary untrusted; no source. Output conservatively labeled as join of all inputs. If binary leaks information, output still restricted by most-restrictive-input label. Conservative labeling compensates for lack of analysis.

\textbf{Alternatives Rejected}:
\begin{itemize}
\item \textbf{Fixed Programs Only}: Inflexible. Hospital cannot add custom analysis without modifying server.
\item \textbf{User-Uploaded Only (White-Box)}: Requires all code to be source available. Incompatible with legacy systems, vendor binaries.
\item \textbf{Black-Box Only}: No formal analysis. Over-restrictive labels (label creep) limit research uses.
\end{itemize}

Hybrid supports both research innovation and legacy system integration.

\subsubsection{Information Flow Method: Denning vs. Volpano vs. Myers-Liskov}

\textbf{Our Choice}: Myers-Liskov Decentralized Model.

\textbf{Why Each Alternative Insufficient}:

\begin{itemize}
\item \textbf{Denning \& Denning (1977)}: First general framework for certification-based information flow. Manual proofs of non-interference. No principled declassification mechanism. Does not support decentralized multi-stakeholder policies.

\item \textbf{Volpano (1996)}: Type-based information flow. Fixed centralized lattice (Low, High). Type system enforces implicit and explicit flow control. Sound and complete for fixed lattices. But:
  \begin{itemize}
  \item Fixed lattice cannot express per-principal policies (patient wants one policy, researcher wants different policy)
  \item No declassification support (or declassification treated as ad-hoc exception)
  \item Doesn't scale to multi-principal systems (thousands of patients with heterogeneous policies)
  \end{itemize}

\item \textbf{Myers-Liskov (1997)}: Decentralized model. Each principal defines independent policies. Labels have form $\text{principal}: \{\text{readers}\}$. Supports:
  \begin{itemize}
  \item \textbf{Per-principal Policies}: Each principal controls own labels independently
  \item \textbf{Label Polymorphism}: Labels can vary at different program points
  \item \textbf{Owner-Authorized Declassification}: Principal $p$ can downgrade labels on $p$-labeled data
  \item \textbf{Multi-Principal Reasoning}: Explicit handling of policies from multiple principals
  \item \textbf{Scalability}: Arbitrary number of principals with heterogeneous policies
  \end{itemize}
\end{itemize}

Myers-Liskov directly addresses hospital scenario (multi-principal, heterogeneous policies, declassification). Volpano and Denning do not.

\subsubsection{Scope: Confidentiality Only vs. Confidentiality + Integrity}

\textbf{Our Choice}: Both confidentiality and integrity.

\textbf{Why}: Hospital cares about both:
\begin{itemize}
\item \textbf{Confidentiality}: Patients' medical data must not leak to unauthorized parties (insurance, employers, public). Ensures privacy.
\item \textbf{Integrity}: Medication records must not be altered by untrusted programs or malicious users. Ensures accurate medical treatment.
\end{itemize}

Information flow labels protect both:
\begin{itemize}
\item Confidentiality labels: $\text{patient}: \{\text{authorized\_readers}\}$. Restricts who can read patient data.
\item Integrity labels: $\text{hospital}: \{\text{authorized\_writers}\}$. Restricts who can modify hospital records (medication dosages, treatment plans).
\end{itemize}

Example: Medication database labeled $\text{hospital}: \{\text{hospital}\}$ (only hospital staff can write). Untrusted program can read but not modify medication instructions.

\textbf{Alternative Rejected}: Confidentiality only. Simpler but incomplete. Hospital also needs guarantee that untrusted programs cannot corrupt medication records or treatment plans.

\subsection{Detailed Design Achievements}

\subsubsection{1. Decentralized Authority with Multi-Principal Policies}

Each principal (patient, doctor, researcher, hospital, analyzer, server) independently defines policies:
\begin{itemize}
\item Patient $P_1$ specifies: data readable only by doctor and patient
\item Hospital specifies: reference database readable only by hospital
\item Researcher specifies: study data readable by research team
\item Analyzer specifies: processed results readable by authorized parties after type-check
\end{itemize}

Contrast with centralized systems: No single administrator controls all policies. Patients maintain autonomy over their data. Policies expressed as labels, composed via join (intersection of readers).

\subsubsection{2. Fine-Grained Access Control via Label Ordering}

Label ordering $L_1 \leq L_2$ encodes access policies. If $L_1$ (more restrictive, fewer readers) can flow to $L_2$ (less restrictive, more readers), information flow is safe.

Example: $\text{patient}_1: \{\} \leq \text{patient}_1: \{\text{doctor}_1\}$ (restrictive patient label can flow to less-restrictive label including doctor). Justifies: doctor can see information derived from patient data.

\subsubsection{3. Secure Program Execution via Type Soundness}

White-box programs type-checked before execution. Three rules guaranteed:
\begin{enumerate}
\item Explicit flows respect label ordering
\item Implicit flows tracked via program counter
\item Declassifications only by authorized principals
\end{enumerate}

Type-check pass $\Rightarrow$ program certified secure. No runtime enforcement needed; static proof sufficient. Prevents information leaks even from malicious or buggy code.

Black-box programs: output conservatively labeled by join of inputs. No source analysis but guaranteed secure by conservative labeling.

\subsubsection{4. Principled Declassification with Audit Trails}

Declassification not ad-hoc. Governed by formal conditions:
\begin{itemize}
\item K-anonymity: $k=5$ ensures re-id risk $\leq 20\%$
\item Aggregation: statistics computed from $\geq k$ individuals, non-recoverable
\item Differential privacy: query responses provably indistinguishable
\end{itemize}

All declassifications logged: timestamp, principal, justification, signature. Enables post-hoc accountability and breach investigation.

\subsubsection{5. Hybrid Confidentiality and Integrity}

Labels support both:
\begin{itemize}
\item Read policies (confidentiality): who can read?
\item Write policies (integrity): who can write?
\end{itemize}

Examples: 
\begin{itemize}
\item Patient data: readable by patient + doctor (confidentiality). Writable by patient only (integrity).
\item Hospital medication records: readable by hospital + doctors (confidentiality). Writable by doctors only (integrity).
\end{itemize}

Information flow analysis ensures untrusted programs cannot violate either policy.

\subsection{Comprehensive Comparison with Volpano Type System}

\begin{table}[h]
\centering
\begin{tabular}{|l|l|l|}
\hline
\textbf{Dimension} & \textbf{Volpano (1996)} & \textbf{Myers-Liskov (Our Design)} \\
\hline
\textbf{Lattice} & Fixed, global & Decentralized, per-principal \\
\hline
\textbf{Labels} & Low, High, etc. & principal:\{readers\} \\
\hline
\textbf{Principals} & Single security admin & Multiple independent principals \\
\hline
\textbf{Declassification} & Not principled & Owner-authorized, auditable \\
\hline
\textbf{Multi-Stakeholder} & No & Yes, via label composition \\
\hline
\textbf{Scalability} & Limited (fixed levels) & Arbitrary principals \\
\hline
\textbf{Hospital Scenario} & Inflexible & Direct fit \\
\hline
\textbf{Research Support} & No & Yes, via declassification workflows \\
\hline
\textbf{Integrity} & Yes (implicitly) & Yes (explicit read+write labels) \\
\hline
\end{tabular}
\end{table}

\subsection{Formal Security Guarantees}

All API commands maintain four core properties:

\begin{enumerate}
\item \textbf{Confidentiality}: Information observable only to authorized principals in label reader set. Unauthorized read attempts return error (no data leak). Type checking ensures no implicit leaks through control flow or timing.

\item \textbf{Non-Interference}: Unauthorized principals' observations (errors, labeled outputs) independent of secret data values. Cannot infer sensitive information from program behavior, output labels, or execution time. Implicit flow control via program counter prevents timing-based inference.

\item \textbf{Type Soundness (White-Box)}: Type-checked programs certified secure before execution. Soundness theorem (Theorem 1, Section 5) guarantees explicit and implicit flow constraints satisfied. Successful type-check $\Rightarrow$ no covert channels in execution.

\item \textbf{Declassification Integrity}: Label changes only by owner principals. Reader expansion monotonic ($R' \supseteq R$). All changes logged with timestamp, justification, digital signature. Audit trail enables accountability and breach investigation.
\end{enumerate}

These guarantees hold even if:
\begin{itemize}
\item Users are dishonest (claim unauthorized access)
\item Programs contain bugs (unintended information flows)
\item Programs are malicious (deliberate data exfiltration attempts)
\item Network is monitored (attacker observes encrypted traffic)
\end{itemize}

Server design ensures no violation of security policy.

\subsection{Limitations and Future Directions}

\subsubsection{1. Label Creep in Black-Box Programs}

Black-box binaries conservatively label output as join of all inputs. With many diverse inputs, reader set shrinks (intersection of readers); output may become over-restrictive.

Example: Output combines patient data (readers: {doctor}) + hospital data (readers: {hospital}) + public data (readers: {public}). Join: readers = {} (empty). Output readable by nobody.

Impact: Limits legitimate uses of black-box program output.

Mitigation: Use white-box analysis when possible. For black-box, designer can annotate intended output label to override conservative join.

\subsubsection{2. Covert Channels Beyond Storage}

Model addresses storage and explicit control-flow information channels. Does not prevent:
\begin{itemize}
\item Timing side-channels (program execution time may correlate with secret values)
\item Cache side-channels (cache hit/miss patterns leak control flow)
\item Power analysis side-channels (power consumption correlates with computation)
\end{itemize}

Mitigating these requires additional mechanisms (constant-time algorithms, cache partitioning, power analysis-resistant hardware).

\subsubsection{3. Label Inference Complexity}

Myers-Liskov paper discusses label inference algorithms for unannotated programs. Full implementation requires sophisticated static analysis to automatically assign security labels to all variables.

Current design assumes programmer-supplied labels. Automatic inference would improve usability but requires sophisticated program analysis.

\subsubsection{4. Declassification Condition Enforcement}

Model supports owner-authorized declassification. Enforcing declassification conditions (k-anonymity, aggregation, etc.) requires additional audit checks beyond label operations.

Example: Analyzer claims "output satisfies k-anonymity." Requires server to verify k-anonymity actually achieved (not just trust analyzer's claim).

Design assumes operators verify conditions. Full automation would require integrating privacy verification tools (k-anonymity checkers, differential privacy analyzers).

\subsubsection{5. Multi-Principal Collaboration on Shared Data}

Current model: each principal controls own labels independently. Doesn't explicitly address scenarios where multiple principals must jointly authorize access.

Example: Two patients' data combined for research. Both patients must consent to sharing.

Workaround: Label uses reader sets intersection (join operation). Both patients' consent required implicitly in label composition.

Formal extension: Explicit multi-principal authorization policies (requires consensus of $k$-of-$n$ principals). Beyond current scope.

\subsection{Summary: Design Completeness and Security}

This design provides:

\begin{itemize}
\item ✓ Complete API specification (5 commands with formal semantics)
\item ✓ Formal security arguments for each command (Sections 5)
\item ✓ Hospital medical scenario fully worked out (3-step declassification, Section 6)
\item ✓ Hybrid confidentiality and integrity protection
\item ✓ Type soundness theorem with proof sketch (Section 5)
\item ✓ Audit trails for accountability and breach investigation
\item ✓ Support for multi-stakeholder collaboration (patients, doctors, researchers, hospital)
\item ✓ Decentralized principal authority (no central bottleneck)
\item ✓ Owner-authorized declassification with formal conditions
\item ✓ 100\% Myers-Liskov framework (no ad-hoc extensions)
\end{itemize}

The system provides formal non-interference guarantees with high practical value for medical data hosting, research collaboration, and multi-stakeholder information sharing scenarios. While limitations (label creep, covert channels, inference complexity) exist, core security properties remain provable against dishonest users and malicious programs.
